\makeabstract
    {
    Development of an Efficient Synthetic Route to Azangulene for the Investigation of Ferroelectric Properties
    }
    {
    Allen Lee\textsuperscript{1}, Valeria Badra\textsuperscript{1}, Ren A. Wiscons\textsuperscript{1}
    }
    {
    \textsuperscript{1} Department of Chemistry, Amherst College, Amherst, MA
    }
    {
    We aimed to develop an efficient synthetic route to azangulene, a nitrogen-containing helicene to investigate ferroelectric properties.  A previously developed synthetic route had been investigated by the laboratory for several years but consistently encountered difficulties in the mid stages of the synthesis.  Further investigation identified a fundamental reactivity conflict: conditions that promoted a high-yielding reaction also promoted undesired debromination, preventing the necessary functional groups from coexisting under the same reaction conditions.  Therefore, a modified synthetic route was developed to overcome this limitation.
    }
    {
    The revised route combines elements of the previous azangulene synthesis with strategies from the reported synthesis of borangulene.  The first step uses a nucleophilic aromatic substitution under basic conditions to construct the extended aromatic precursor.  The resulting nitro compound is then pushed through a reduction using hydrazine hydrate and Pd/C in ethanol at \ensuremath{>}100 {\textdegree}C.  Reaction products were monitored and purified using standard organic synthesis techniques.  
    }
    {
    Investigation of the previous synthetic route demonstrated that high yield could not be achieved without high probability of debromination, making the original strategy unsuitable for reliably producing the desired intermediate.  An alternative route was therefore developed by combining the previous azangulene strategy with elements of the borangulene synthesis.  The first two steps of the revised route have been successfully carried out.  Including formation of the nitro precursor and its subsequent reduction to the corresponding amine.  However, the reduction step is still an area requiring further optimization to obtain consistently high yields with less time spent.  
    }
    {
    Identification of the competing debromination pathway provided an explanation from the limitations of the previous azangulene synthesis and motivated the development of a revised synthetic route.  The successful completion of the first two steps demonstrates initial progress toward the new route.  Future work will focus on optimizing the reduction step and determining suitable conditions for the remaining steps.  Successful completion of the revised synthesis would provide a more efficient route to azangulene for the following investigation of its ferroelectric properties.  
    }

    \newpage
    